\section{Compilación, programación y cambio de constantes}
\label{sec:operacion}

Todas las herramientas vienen en OSS CAD Suite \cite{osscad}: Yosys, nextpnr, IceStorm,
Icarus Verilog y Python 3. Desde su terminal, en la carpeta del proyecto:

\begin{table}[H]
\centering
\caption{Comandos del Makefile.}
\label{tab:make}
\begin{tabularx}{\linewidth}{l X}
\toprule
Comando & Qué hace \\
\midrule
\code{make prog} & ensambla \code{game.s}, sintetiza, hace place and route, genera \code{game\_top.bin} y programa la Go Board \\
\code{make} & lo mismo, sin programar \\
\code{make hex} & solo \code{sw/game.s} $\to$ \code{sw/game.hex} \\
\code{make test} & pruebas del assembler \\
\code{make sim} & prueba de instrucciones y simulación del juego \\
\bottomrule
\end{tabularx}
\end{table}

El Makefile conoce las dependencias: \code{game\_top.bin} depende de \code{game.hex},
que depende de \code{game.s} y del assembler. Por eso, tras editar \code{game.s}, basta
\code{make prog}: regenera solo lo necesario. La compilación completa toma unos
8 segundos.

Para regenerar \code{game.hex} sin make:

\begin{lstlisting}[language=bash]
python3 asm/pochoasm.py sw/game.s -o sw/game.hex
\end{lstlisting}

\paragraph{Cambiar una constante.} Durante la demostración se pedirá modificar una
constante del juego. Todas están con nombre al comienzo de \code{sw/game.s}
(código~\ref{lst:constantes}). Por ejemplo, para jugar cinco rondas en vez de diez:

\begin{enumerate}
  \item Editar la línea \code{.equ ROUNDS, 10} a \code{.equ ROUNDS, 5}.
  \item Ejecutar \code{make prog}.
\end{enumerate}

El promedio se sigue calculando bien porque divide por \code{ROUNDS}, no por 10. De la
misma forma, \code{.equ WAIT\_CYCLES, 125000000} deja la fase de los cuatro LEDs en
5 segundos. El \code{.hex} cambia porque la RAM se inicializa al sintetizar, así que un
cambio en el programa exige volver a sintetizar el diseño completo.
